\section{问题三：带时间要求的多日排班}

\subsection{问题分析}

\subsubsection{问题背景与难点}

问题三在问题二基础上叠加时间窗、机队数量（24 架）和运营时段（06:00--20:00）三层约束。与问题一、二相比，问题三的核心差异在于：

\begin{itemize}
\item \textbf{时间窗约束}：每条需求具有最早离开起点时刻和最晚送达时刻，需求时间在 2026-08-03 至 2026-08-10 之间
\item \textbf{机队约束}：飞机数量有限（24 架），需考虑机场-机型组合的可用飞机数
\item \textbf{运营时段约束}：飞机 06:00--18:00 起飞且不晚于 20:00 返回，不得在海上设施过夜
\item \textbf{临时任务}：阶段二需在不增加总飞机使用时间的前提下插入 160 人临时任务
\end{itemize}

\subsubsection{核心矛盾与数据洞察}

对 \texttt{peopleQ3.csv} 的统计揭示关键数据特征：84.1\% 的需求时间窗跨天（均值 41.6 小时），为需求-日期分配提供充分自由度；但时间窗碎片化使单架次平均仅服务 7--17 人（座位 12--19），大量座位被浪费。核心矛盾是\textbf{时间窗碎片化与运力效率的冲突}——利用率从问题二的 88.18\% 骤降至 35.81\%（见图~\ref{fig:demand_dist} 的需求时间分布）。

\subsection{模型建立}

\subsubsection{问题分解与预处理}

问题三的规模是问题二的 7 倍（7 天 $\times$ 4000 人 $\times$ 24 架飞机），直接求解不可行。但从题目约束中可提炼出关键洞察，实现严格等价的降维分解。

\begin{theorem}
\label{thm:decouple}
给定约束`飞机不得在海上设施过夜、每个架次必须返回所属机场、飞机 06:00--18:00 起飞且不晚于 20:00 返回''，任意架次的起飞时刻与返回时刻必在同一自然日内。
\end{theorem}

\begin{proof}
设架次在 $ 日 $ 时刻起飞（ \geq 06{:}00$）。若返回时刻落在 +1$ 日，则飞机在海上设施停留跨过午夜，违反`不得在海上设施过夜''；若飞机在返回前停留于机场，则架次已在 $ 日结束。故架次不可能跨天。
\end{proof}

\begin{corollary}
\label{cor:indep}
由于架次不跨天、飞机每日从所属机场重置出发，各日之间不存在状态转移耦合。多日排班问题可精确分解为  = 7$ 个独立的单日问题：
$\min \sum_{d=1}^{7} Z_d$，其中 $ 为第 $ 日最优目标值。
\end{corollary}

\begin{corollary}
\label{cor:dateset}
对需求 $，其可行日期集合为  = \{ d : [e_p, l_p] \cap [06{:}00, 20{:}00]_d \neq \emptyset \}$。由于 84.1\% 的需求时间窗跨天，$|D_p| \geq 2$ 常成立，为需求-日期分配提供自由度。
\end{corollary}

在此基础上，进一步给出时间窗的预处理方法。

\begin{proposition}
\label{prop:cpm}
对候选路径  = [A, S_1, \ldots, S_K, A]$，定义前向最早到达时刻与后向最晚离开时刻：
\begin{equation}
\text{EAT}_k = \max\!\left(\text{EAT}_{k-1} + \frac{d_{k-1,k}}{v} \times 60 + s_{k-1},\; \max_{p \in P_k^{\mathrm{b}}} e_p\right)
\label{eq:eat}
\end{equation}
\begin{equation}
\text{LDT}_k = \min\!\left(\text{LDT}_{k+1} - \frac{d_{k,k+1}}{v} \times 60,\; \min_{p \in P_k^{\mathrm{a}}} l_p\right)
\label{eq:ldt}
\end{equation}

其中 $ 为停靠时间（加油 20 min / 普通 10 min），$ 为机型速度，边界条件 $\text{EAT}_0 = 06{:}00$，$\text{LDT}_{K+1} = 20{:}00$。若 $\exists\, k$ 使得 $\text{EAT}_k > \text{LDT}_k$，则路径 $ 不可行，直接剔除。此收缩为纯数学预处理，以 (K)$ 代价在路径生成阶段剔除时间窗不可行路径。
\end{proposition}

此外，按天解耦的贪心分配跨天需求存在可量化的损失上界：设跨天需求占比为 $\alpha$，单日 MIP 相对间隙为 $\delta$，单条跨天需求分配到非最优日期的额外成本上界为 $\Delta$，则 $\hat{T} \leq (1+\delta) T^* + \alpha |\mathcal{P}| \Delta$。本数据 $\alpha = 84.1\%$，MIP 间隙是主要瓶颈而非解耦损失。
